% ============================================================
\chapter{Conditional Probability and Independence}
\label{ch:conditional}
% ============================================================

Imagine you are a doctor. A patient tests positive for a rare disease. The test is 99\% accurate. Should you be worried? Most people---including many physicians---would say yes without hesitation. And yet, if the disease affects only one in ten thousand people, the patient is far more likely to be healthy than sick. This unsettling result, which has fooled experts for centuries, is a direct consequence of conditional probability.

The idea itself traces back to Thomas Bayes, a Presbyterian minister whose posthumous essay, published in 1763 by his friend Richard Price, laid the groundwork for reasoning under uncertainty. But it was Pierre-Simon Laplace who, independently and with far greater mathematical ambition, turned conditional probability into a systematic tool. The story of this chapter is the story of how a simple question---``what changes when we learn something new?''---reshapes the entire landscape of probability theory.

\begin{intuition}
Conditional probability allows us to update our information when an event is
observed. Independence formalises the idea that two events have no influence on
each other. These two concepts lie at the heart of all probability theory.
\end{intuition}

% ------------------------------------------------------------
\section{Conditional Probability}
% ------------------------------------------------------------

\begin{definition}[Conditional probability]
Let $(\Omega, \mathcal{F}, \PP)$ be a probability space and $B \in \mathcal{F}$
an event with $\PP(B) > 0$. The \textbf{conditional probability} of $A$ given $B$
is
\[
    \PP(A \mid B) = \frac{\PP(A \cap B)}{\PP(B)}.
\]
\end{definition}

\begin{proposition}
For fixed $B$ with $\PP(B) > 0$, the map $A \mapsto \PP(A \mid B)$ is a
probability measure on $(\Omega, \mathcal{F})$.
\end{proposition}

\begin{proof}
We verify the Kolmogorov axioms:
\begin{itemize}
    \item $\PP(A \mid B) = \PP(A \cap B)/\PP(B) \geq 0$ for all $A$.
    \item $\PP(\Omega \mid B) = \PP(B)/\PP(B) = 1$.
    \item If $(A_n)$ are mutually disjoint, then so are $(A_n \cap B)$, and by
    $\sigma$-additivity of $\PP$:
    \[
        \PP\!\left(\bigcup_n A_n \;\middle|\; B\right)
        = \frac{\PP\bigl(\bigcup_n (A_n \cap B)\bigr)}{\PP(B)}
        = \frac{\sum_n \PP(A_n \cap B)}{\PP(B)}
        = \sum_n \PP(A_n \mid B). \qedhere
    \]
\end{itemize}
\end{proof}

\begin{example}
Two fair dice are rolled. Given that the sum is~8, what is the probability that
the first die shows~3?

Let $A = \{\text{first die} = 3\}$ and $B = \{\text{sum} = 8\}$.
The outcomes giving sum~8 are $(2,6),(3,5),(4,4),(5,3),(6,2)$,
so $\PP(B) = 5/36$. Since $A \cap B = \{(3,5)\}$, $\PP(A \cap B) = 1/36$.
Hence $\PP(A \mid B) = (1/36)/(5/36) = 1/5$.
\end{example}

% ------------------------------------------------------------
\section{Chain Rule and Bayes' Theorem}
% ------------------------------------------------------------

\begin{theorem}[Multiplication rule / chain rule]
For events $A_1, \ldots, A_n$ with $\PP(A_1 \cap \cdots \cap A_{n-1}) > 0$:
\[
    \PP(A_1 \cap \cdots \cap A_n)
    = \PP(A_1)\,\PP(A_2 \mid A_1)\,\PP(A_3 \mid A_1 \cap A_2)
    \cdots \PP(A_n \mid A_1 \cap \cdots \cap A_{n-1}).
\]
\end{theorem}

\begin{definition}[Partition of the sample space]
A family $(B_i)_{i \in I}$ of events forms a \textbf{partition} of $\Omega$ if:
\begin{enumerate}[label=(\roman*)]
    \item $B_i \cap B_j = \emptyset$ for $i \neq j$ (disjoint);
    \item $\bigcup_{i \in I} B_i = \Omega$ (exhaustive);
    \item $\PP(B_i) > 0$ for all $i$.
\end{enumerate}
\end{definition}

\begin{theorem}[Law of total probability]
\label{thm:total_prob}
Let $(B_i)_{i \in I}$ be a partition of $\Omega$ (with $I$ finite or countable).
For any $A \in \mathcal{F}$:
\[
    \PP(A) = \sum_{i \in I} \PP(A \mid B_i)\,\PP(B_i).
\]
\end{theorem}

\begin{proof}
Since $A = A \cap \Omega = \bigcup_i (A \cap B_i)$ (disjoint union):
\[
    \PP(A) = \sum_i \PP(A \cap B_i) = \sum_i \PP(A \mid B_i)\,\PP(B_i). \qedhere
\]
\end{proof}

\begin{theorem}[Bayes' theorem]
\label{thm:bayes}
Let $(B_i)_{i \in I}$ be a partition of $\Omega$ and $A$ an event with
$\PP(A) > 0$. Then for every $j \in I$:
\[
    \PP(B_j \mid A) = \frac{\PP(A \mid B_j)\,\PP(B_j)}%
    {\sum_{i \in I} \PP(A \mid B_i)\,\PP(B_i)}.
\]
\end{theorem}

\begin{intuition}[title=\textbf{Bayesian interpretation}]
The $\PP(B_i)$ are the \emph{prior} probabilities of the hypotheses $B_i$.
After observing $A$, Bayes' theorem yields the \emph{posterior} probabilities
$\PP(B_i \mid A)$, updated via the likelihood $\PP(A \mid B_i)$.
\end{intuition}

\begin{example}[Medical screening]
A test has sensitivity 99\% (true positive) and specificity 95\% (true negative).
The disease prevalence is 0.1\%. What is the probability of disease given a
positive test?

Let $D$ = ``diseased'' and $T^+$ = ``test positive''.
$\PP(D) = 0.001$, $\PP(T^+ \mid D) = 0.99$, $\PP(T^+ \mid D^c) = 0.05$.
\begin{align*}
    \PP(D \mid T^+)
    &= \frac{\PP(T^+ \mid D)\,\PP(D)}%
    {\PP(T^+ \mid D)\,\PP(D) + \PP(T^+ \mid D^c)\,\PP(D^c)} \\
    &= \frac{0.99 \times 0.001}{0.99 \times 0.001 + 0.05 \times 0.999}
    = \frac{0.00099}{0.05094}
    \approx 0.0194.
\end{align*}
Even with an excellent test, the probability of disease given a positive result is
only about 2\%, because prevalence is low.
\end{example}

% ------------------------------------------------------------
\section{Independence of Events}
% ------------------------------------------------------------

\begin{definition}[Independence of two events]
Two events $A$ and $B$ are \textbf{independent}, written $A \indep B$, if
\[
    \PP(A \cap B) = \PP(A) \cdot \PP(B).
\]
\end{definition}

\begin{remark}
If $\PP(B) > 0$, independence is equivalent to $\PP(A \mid B) = \PP(A)$:
knowing that $B$ occurred does not change the probability of $A$.
\end{remark}

\begin{warning}[title=\textbf{Independence $\neq$ disjointness}]
Two disjoint events ($A \cap B = \emptyset$) with positive probabilities are
\emph{never} independent, since $\PP(A \cap B) = 0 \neq \PP(A)\PP(B)$.
Do not confuse these notions!
\end{warning}

\begin{definition}[Mutual independence]
Events $A_1, \ldots, A_n$ are \textbf{mutually independent} if, for every
subset $J \subseteq \{1, \ldots, n\}$ with $\abs{J} \geq 2$:
\[
    \PP\!\left(\bigcap_{j \in J} A_j\right) = \prod_{j \in J} \PP(A_j).
\]
\end{definition}

\begin{warning}[title=\textbf{Pairwise $\neq$ mutual independence}]
Pairwise independence does not imply mutual independence. All sub-families must
be checked, not just pairs.
\end{warning}

\begin{example}[Bernstein's counterexample]
Roll two fair dice. Define:
\begin{itemize}
    \item $A$ = ``first die is even'',
    \item $B$ = ``second die is even'',
    \item $C$ = ``the sum is even''.
\end{itemize}
One verifies $\PP(A) = \PP(B) = \PP(C) = 1/2$ and every pair is independent:
$\PP(A \cap B) = 1/4 = \PP(A)\PP(B)$, etc. However,
$\PP(A \cap B \cap C) = 1/4 \neq 1/8 = \PP(A)\PP(B)\PP(C)$.
The three events are pairwise but not mutually independent.
\end{example}

\begin{proposition}
If $A$ and $B$ are independent, then so are $(A, B^c)$, $(A^c, B)$, and
$(A^c, B^c)$.
\end{proposition}

\begin{proof}
$\PP(A \cap B^c) = \PP(A) - \PP(A \cap B) = \PP(A) - \PP(A)\PP(B)
= \PP(A)(1 - \PP(B)) = \PP(A)\PP(B^c)$.
The other cases are analogous.
\end{proof}

% ------------------------------------------------------------
\section{The Second Borel--Cantelli Lemma}
% ------------------------------------------------------------

\begin{theorem}[Second Borel--Cantelli lemma]
\label{thm:borel_cantelli_2}
Let $(A_n)_{n \geq 1}$ be a sequence of \textbf{mutually independent} events.
If $\sum_{n=1}^{\infty} \PP(A_n) = +\infty$, then
\[
    \PP\!\left(\limsup_{n \to \infty} A_n\right) = 1.
\]
That is, almost surely infinitely many of the $A_n$ occur.
\end{theorem}

\begin{proof}
It suffices to show $\PP\bigl(\bigcup_{k=n}^{\infty} A_k\bigr) = 1$ for all $n$.
Taking complements:
\[
    \PP\!\left(\bigcap_{k=n}^{N} A_k^c\right)
    = \prod_{k=n}^{N} \PP(A_k^c)
    = \prod_{k=n}^{N} (1 - \PP(A_k))
    \leq \prod_{k=n}^{N} e^{-\PP(A_k)}
    = \exp\!\left(-\sum_{k=n}^{N} \PP(A_k)\right).
\]
As $N \to \infty$, the sum diverges, so the product tends to~0. By continuity
from above: $\PP\bigl(\bigcap_{k=n}^{\infty} A_k^c\bigr) = 0$, hence
$\PP\bigl(\bigcup_{k=n}^{\infty} A_k\bigr) = 1$.
\end{proof}

\begin{keyformulas}[title=\textbf{Summary: Borel--Cantelli}]
\begin{itemize}
    \item $\sum \PP(A_n) < \infty \implies \PP(\limsup A_n) = 0$ \quad (no hypothesis needed).
    \item $\sum \PP(A_n) = \infty$ \textbf{and} independence
          $\implies \PP(\limsup A_n) = 1$.
\end{itemize}
\end{keyformulas}

% ------------------------------------------------------------
\section{Independence of $\sigma$-Algebras and Random Variables}
% ------------------------------------------------------------

\begin{definition}[Independent $\sigma$-algebras]
Two sub-$\sigma$-algebras $\mathcal{G}_1, \mathcal{G}_2 \subseteq \mathcal{F}$
are \textbf{independent} if, for all $A \in \mathcal{G}_1$ and
$B \in \mathcal{G}_2$:
\[
    \PP(A \cap B) = \PP(A)\,\PP(B).
\]
\end{definition}

\begin{definition}[Independent random variables --- preview]
Two random variables $X$ and $Y$ are \textbf{independent} if the $\sigma$-algebras
they generate, $\sigma(X)$ and $\sigma(Y)$, are independent. In terms of
distributions (see Chapter~3 for full definitions), this amounts to:
\[
    \PP(X \in A,\; Y \in B) = \PP(X \in A)\,\PP(Y \in B)
\]
for all Borel sets $A, B \in \mathcal{B}(\R)$.
\end{definition}

% ------------------------------------------------------------
\section{Classical Applications}
% ------------------------------------------------------------

\begin{example}[The Monty Hall problem]
A contestant chooses one of three doors. Behind one is a car; behind the other
two are goats. The host, who knows where the car is, opens a door revealing a
goat, then offers the contestant the chance to switch.

Let $C_i$ = ``the car is behind door $i$'' ($i=1,2,3$), and suppose the contestant
picks door~1. By Bayes' theorem, the probability of winning by switching is $2/3$,
while staying gives $1/3$. Switching is therefore advantageous.
\end{example}

\begin{example}[Gambler's ruin]
A gambler starts with capital $k \in \{1, \ldots, N-1\}$. At each round, they win
\$1 with probability $p$ or lose \$1 with probability $q = 1-p$, independently.
Play stops when the capital reaches 0 (ruin) or $N$ (target).

Using conditional probability and a recurrence argument, the ruin probability is:
\[
    \PP(\text{ruin} \mid X_0 = k) =
    \begin{cases}
        \dfrac{(q/p)^k - (q/p)^N}{1 - (q/p)^N} & \text{if } p \neq q, \\[6pt]
        1 - \dfrac{k}{N} & \text{if } p = q = \tfrac{1}{2}.
    \end{cases}
\]
\end{example}

% ------------------------------------------------------------
\section{Exercises}
% ------------------------------------------------------------

\begin{exercise}
A batch contains 10 items, 3 of which are defective. Two items are drawn
successively without replacement. Find the probability that the second is
defective given that the first is defective.
\end{exercise}

\begin{exercise}
Three machines $M_1, M_2, M_3$ produce 50\%, 30\%, and 20\% of a factory's
output respectively. Their defect rates are 2\%, 3\%, and 5\%. A randomly chosen
item is defective. What is the probability it came from $M_2$?
\end{exercise}

\begin{exercise}
Show that if $A$, $B$, $C$ are mutually independent, then $A \cup B$ and $C$
are independent.
\end{exercise}

\begin{exercise}
A fair coin is tossed $n$ times. Let $A_k$ be the event ``the $k$-th toss is
heads''. Show that the $A_k$ are mutually independent.
\end{exercise}

\begin{exercise}[Iterated Bayes]
An urn contains 1 red and 1 blue ball. A ball is drawn at random, replaced, and
a ball of the same colour is added. After $n$ draws, find the probability that
all drawn balls were red.
\end{exercise}

\begin{exercise}
Let $(A_n)_{n \geq 1}$ be independent events with $\PP(A_n) = 1/(n+1)$.
Show that $\PP(\limsup A_n) = 1$. Interpret the result.
\end{exercise}

\begin{exercise}[Simpson's paradox]
Give a numerical example where $\PP(A \mid B) > \PP(A \mid B^c)$ but
$\PP(A \mid B \cap C) < \PP(A \mid B^c \cap C)$ and
$\PP(A \mid B \cap C^c) < \PP(A \mid B^c \cap C^c)$. Explain the phenomenon.
\end{exercise}

\begin{exercise}
Let $p \in (0,1)$ and $(X_n)_{n \geq 1}$ be independent Bernoulli trials with
parameter $p$. Show that $\PP(\exists\, n : X_n = 1) = 1$.
\end{exercise}
